\documentclass[11pt]{article}

% ----------------------------
% Packages
% ----------------------------
\usepackage[a4paper,margin=1in]{geometry}
\usepackage{hyperref}
\usepackage{listings}
\usepackage{xcolor}
\usepackage{enumitem}
\usepackage{titlesec}
\usepackage{booktabs}
\usepackage{amsmath}

% ----------------------------
% Styling
% ----------------------------
\titleformat{\section}{\large\bfseries}{\thesection}{1em}{}
\titleformat{\subsection}{\normalsize\bfseries}{\thesubsection}{1em}{}
\titleformat{\subsubsection}{\normalsize\bfseries}{\thesubsubsection}{1em}{}

\definecolor{codegray}{rgb}{0.95,0.95,0.95}
\definecolor{commentgreen}{rgb}{0.0,0.5,0.0}

\lstset{
  basicstyle=\ttfamily\small,
  backgroundcolor=\color{codegray},
  frame=single,
  breaklines=true,
  showstringspaces=false,
  keywordstyle=\color{blue},
  commentstyle=\color{commentgreen},
  stringstyle=\color{red}
}

% ----------------------------
% Document Info
% ----------------------------
\title{\textbf{MCV Lab Standard: Scalable ML Experimentation}\\
\large Reproducibility, Config Management, and HPC Execution (Argparse + Hydra + SLURM)}
\author{Medical Computer Vision Lab (MCV Lab)}
\date{\today}

\begin{document}
\maketitle

\section{Purpose and Scope}
This document defines the MCV Lab standard for running machine learning experiments in a scalable, reproducible, and collaborative way.
It applies to:
\begin{itemize}[nosep]
    \item all training scripts (vision, time-series, multimodal)
    \item all experiment pipelines (pretrain, finetune, evaluate, inference)
    \item all compute platforms (local GPU workstation, lab server, HPC/SLURM)
\end{itemize}

\section{Core Principles}
\subsection{Non-negotiables}
\begin{enumerate}[label=\textbf{P\arabic*.}, leftmargin=*]
    \item \textbf{No code edits for experiments.}\\
    Experiments must be configured through CLI arguments and/or config overrides only.
    \item \textbf{Full reproducibility.}\\
    Every run must be repeatable from a stored configuration.
    \item \textbf{Traceability.}\\
    Every run must save: configuration, launch command, outputs, checkpoints.
    \item \textbf{Scalability.}\\
    The code must support sweeps, multiple datasets, and automation (bash/SLURM/Hydra multirun).
\end{enumerate}

\subsection{Why this matters}
The lab frequently runs hundreds of experiments (hyperparameter sweeps, ablations, domain shifts, modality missingness).
Without strict standards, results become:
\begin{itemize}[nosep]
    \item not reproducible
    \item hard to interpret
    \item difficult to compare fairly
    \item expensive in wasted compute time
\end{itemize}

\section{Experiment Taxonomy: When to Use Argparse vs Hydra}
\subsection{Recommended default}
\begin{itemize}[nosep]
    \item \textbf{Argparse} for small single-script projects and quick baselines.
    \item \textbf{Hydra} for any project expected to scale (recommended).
\end{itemize}

\subsection{Decision table}
\begin{center}
\begin{tabular}{ll}
\toprule
\textbf{Scenario} & \textbf{Recommended Tool} \\
\midrule
Quick baseline training / single dataset & Argparse \\
Small ablation study (10--20 runs) & Argparse or Hydra \\
Multiple datasets + models + optimizers & Hydra \\
Multimodal fusion experiments & Hydra \\
Continual learning setups & Hydra \\
Large sweeps on SLURM & Hydra + SLURM arrays \\
\bottomrule
\end{tabular}
\end{center}

\section{Standard Run Output Structure}
All experiments must follow:

\begin{lstlisting}
project/
  runs/
    <run_name>/
      config.json OR config_resolved.yaml
      cmd.txt
      checkpoints/
        last.ckpt
        best.ckpt
      logs/
        metrics.jsonl OR tensorboard OR wandb
\end{lstlisting}

\subsection{Minimum required artifacts}
Every run directory must include:
\begin{itemize}[nosep]
    \item \texttt{cmd.txt}: exact command used to launch the experiment
    \item resolved config (\texttt{config.json} or \texttt{config\_resolved.yaml})
    \item at least one checkpoint
\end{itemize}

\section{Lab Rule: No Code Changes for Experiments}
\subsection{Allowed}
\begin{itemize}[nosep]
    \item \texttt{python train.py --lr 1e-4 --seed 2}
    \item \texttt{python train.py --config configs/baseline.json --batch\_size 64}
    \item \texttt{python train.py model=resnet50 optim.lr=5e-5} (Hydra)
\end{itemize}

\subsection{Disallowed}
\begin{itemize}[nosep]
    \item editing learning rate/batch size/epochs in \texttt{train.py}
    \item hardcoding dataset paths
    \item commenting/uncommenting blocks to switch models
\end{itemize}

\section{Reproducibility Standard}
\subsection{Required}
All training scripts must support:
\begin{itemize}[nosep]
    \item \texttt{seed} argument
    \item determinism toggle (\texttt{--deterministic} or Hydra field)
    \item saving configuration and command
\end{itemize}

\subsection{Recommended}
\begin{itemize}[nosep]
    \item log git commit hash into run directory
    \item log environment: Python version, torch version, CUDA version
\end{itemize}

\section{Part I: Argparse Standard (Baseline)}
\subsection{Why argument parsers scale experimentation}
Argument parsers enable scalability because they:
\begin{itemize}[nosep]
    \item remove the need for code edits
    \item enable easy sweeps via bash/SLURM
    \item enforce safe typed arguments (avoid invalid runs)
    \item support config-file driven defaults with CLI overrides
\end{itemize}

\subsection{Required arguments (minimum)}
Every argparse training script should include:

\begin{itemize}[nosep]
    \item Experiment: \texttt{--run\_name}, \texttt{--out\_dir}
    \item Dataset: \texttt{--dataset}, \texttt{--data\_dir}, \texttt{--num\_workers}
    \item Model: \texttt{--model}, \texttt{--pretrained}
    \item Optimization: \texttt{--epochs}, \texttt{--batch\_size}, \texttt{--lr}, \texttt{--weight\_decay}
    \item Reproducibility: \texttt{--seed}, \texttt{--deterministic}
    \item Hardware: \texttt{--device}, \texttt{--amp}
    \item Resume: \texttt{--resume}
    \item Config support: \texttt{--config} (mandatory)
\end{itemize}

\subsection{Reference template (argparse)}
Below is the lab reference entrypoint template (students may adapt but preserve principles):

\begin{lstlisting}[language=Python]
#!/usr/bin/env python3
"""
train.py - scalable training entrypoint template
"""

from __future__ import annotations

import argparse
import json
import random
import shutil
import sys
from dataclasses import asdict, dataclass
from datetime import datetime
from pathlib import Path
from typing import Any

import numpy as np

try:
    import torch
except ImportError:
    torch = None


def set_seed(seed: int, deterministic: bool = False) -> None:
    random.seed(seed)
    np.random.seed(seed)
    if torch is not None:
        torch.manual_seed(seed)
        torch.cuda.manual_seed_all(seed)
        if deterministic:
            torch.backends.cudnn.deterministic = True
            torch.backends.cudnn.benchmark = False


@dataclass
class TrainConfig:
    run_name: str
    out_dir: str
    dataset: str
    data_dir: str
    model: str
    epochs: int
    batch_size: int
    lr: float
    weight_decay: float
    seed: int
    deterministic: bool
    device: str
    amp: bool
    resume: str | None


def timestamp() -> str:
    return datetime.now().strftime("%Y%m%d_%H%M%S")


def ensure_dir(path: str | Path) -> Path:
    path = Path(path)
    path.mkdir(parents=True, exist_ok=True)
    return path


def dump_json(path: Path, obj: dict[str, Any]) -> None:
    with open(path, "w", encoding="utf-8") as f:
        json.dump(obj, f, indent=2, sort_keys=True)


def load_json(path: str | Path) -> dict[str, Any]:
    with open(path, "r", encoding="utf-8") as f:
        return json.load(f)


def build_parser() -> argparse.ArgumentParser:
    parser = argparse.ArgumentParser(
        formatter_class=argparse.ArgumentDefaultsHelpFormatter
    )
    parser.add_argument("--run_name", type=str, default="")
    parser.add_argument("--out_dir", type=str, default="runs")

    parser.add_argument("--dataset", type=str, default="dummy",
                        choices=["dummy", "ptbxl", "mimic_ecg", "rsna_trauma"])
    parser.add_argument("--data_dir", type=str, default="data")

    parser.add_argument("--model", type=str, default="resnet18",
                        choices=["resnet18", "resnet50", "vit_b16"])

    parser.add_argument("--epochs", type=int, default=30)
    parser.add_argument("--batch_size", type=int, default=16)
    parser.add_argument("--lr", type=float, default=1e-4)
    parser.add_argument("--weight_decay", type=float, default=1e-4)

    parser.add_argument("--seed", type=int, default=0)
    parser.add_argument("--deterministic", action="store_true")

    parser.add_argument("--device", type=str, default="cuda",
                        choices=["cuda", "cpu", "mps"])
    parser.add_argument("--amp", action="store_true")

    parser.add_argument("--resume", type=str, default=None)

    parser.add_argument("--config", type=str, default=None,
                        help="Optional JSON config to load defaults from.")

    return parser


def parse_args_with_config(parser: argparse.ArgumentParser, argv: list[str]) -> argparse.Namespace:
    partial, _ = parser.parse_known_args(argv)
    if partial.config is not None:
        cfg = load_json(partial.config)
        parser.set_defaults(**cfg)
    return parser.parse_args(argv)


def main(argv: list[str]) -> int:
    parser = build_parser()
    args = parse_args_with_config(parser, argv)

    if args.run_name.strip() == "":
        args.run_name = f"{args.dataset}_{args.model}_lr{args.lr}_{timestamp()}"

    run_dir = ensure_dir(Path(args.out_dir) / args.run_name)
    ensure_dir(run_dir / "checkpoints")

    if args.config is not None:
        shutil.copy(args.config, run_dir / "config_used.json")

    cfg = TrainConfig(
        run_name=args.run_name,
        out_dir=str(run_dir),
        dataset=args.dataset,
        data_dir=args.data_dir,
        model=args.model,
        epochs=args.epochs,
        batch_size=args.batch_size,
        lr=args.lr,
        weight_decay=args.weight_decay,
        seed=args.seed,
        deterministic=args.deterministic,
        device=args.device,
        amp=args.amp,
        resume=args.resume,
    )

    dump_json(run_dir / "config.json", asdict(cfg))
    (run_dir / "cmd.txt").write_text("python " + " ".join(argv) + "\n")

    # TODO: training logic
    print(f"Training would run here. Outputs: {run_dir}")
    return 0


if __name__ == "__main__":
    raise SystemExit(main(sys.argv[1:]))
\end{lstlisting}

\section{Part II: Hydra Standard (Recommended)}
\subsection{Why Hydra scales better than argparse}
Hydra enables:
\begin{itemize}[nosep]
    \item hierarchical configs (dataset/model/optim/trainer)
    \item composable config groups
    \item built-in multi-run sweeps (\texttt{-m})
    \item automatic run directory management and config saving
\end{itemize}

\subsection{Required project structure}
\begin{lstlisting}
project/
  train.py
  configs/
    config.yaml
    dataset/
      rsna_trauma.yaml
      mimic_ecg.yaml
    model/
      resnet50.yaml
      vit_b16.yaml
    optim/
      adamw.yaml
      sgd.yaml
    trainer/
      default.yaml
\end{lstlisting}

\subsection{Hydra entrypoint template}
\begin{lstlisting}[language=Python]
#!/usr/bin/env python3
from __future__ import annotations

import os
import random
from pathlib import Path

import numpy as np
import hydra
from omegaconf import DictConfig, OmegaConf

try:
    import torch
except ImportError:
    torch = None


def set_seed(seed: int, deterministic: bool = False) -> None:
    random.seed(seed)
    np.random.seed(seed)
    if torch is not None:
        torch.manual_seed(seed)
        torch.cuda.manual_seed_all(seed)
        if deterministic:
            torch.backends.cudnn.deterministic = True
            torch.backends.cudnn.benchmark = False


@hydra.main(version_base=None, config_path="configs", config_name="config")
def main(cfg: DictConfig) -> None:
    run_dir = Path(os.getcwd())

    (run_dir / "config_resolved.yaml").write_text(
        OmegaConf.to_yaml(cfg), encoding="utf-8"
    )

    set_seed(int(cfg.seed), deterministic=bool(cfg.deterministic))

    print("Run dir:", run_dir)
    print(OmegaConf.to_yaml(cfg))

    # TODO: build dataset/model/optim/trainer based on cfg groups


if __name__ == "__main__":
    main()
\end{lstlisting}

\subsection{Hydra base config}
\begin{lstlisting}
# configs/config.yaml
defaults:
  - dataset: rsna_trauma
  - model: resnet50
  - optim: adamw
  - trainer: default
  - _self_

seed: 0
deterministic: false
\end{lstlisting}

\subsection{Example Hydra multirun sweeps}
\begin{lstlisting}
python train.py -m optim.lr=1e-3,1e-4,1e-5 seed=0,1,2
python train.py -m model=resnet50,vit_b16 dataset=rsna_trauma,mimic_ecg
\end{lstlisting}

\section{Part III: SLURM Standard (HPC Execution)}
\subsection{Required folder structure}
\begin{lstlisting}
project/
  runs/
  slurm_logs/
  scripts/
    run_train.slurm
    sweep_seeds.slurm
    ddp_train.slurm
\end{lstlisting}

\subsection{Baseline SLURM job (single run)}
\begin{lstlisting}[language=bash]
#!/bin/bash
#SBATCH --job-name=mcv_train
#SBATCH --output=slurm_logs/%x_%j.out
#SBATCH --error=slurm_logs/%x_%j.err
#SBATCH --time=08:00:00
#SBATCH --partition=gpu
#SBATCH --gres=gpu:1
#SBATCH --cpus-per-task=8
#SBATCH --mem=32G

set -euo pipefail

echo "JOB_ID: ${SLURM_JOB_ID}"
echo "NODELIST: ${SLURM_NODELIST}"
echo "START_TIME: $(date)"

cd ${SLURM_SUBMIT_DIR}
nvidia-smi || true

python train.py \
  --dataset rsna_trauma \
  --model resnet50 \
  --lr 1e-4 \
  --seed 0 \
  --amp \
  --out_dir runs \
  --run_name "rsna_resnet50_seed0_lr1e4_${SLURM_JOB_ID}"

echo "END_TIME: $(date)"
\end{lstlisting}

\subsection{SLURM arrays for sweeps (recommended)}
\begin{lstlisting}[language=bash]
#!/bin/bash
#SBATCH --job-name=mcv_seed_sweep
#SBATCH --output=slurm_logs/%x_%A_%a.out
#SBATCH --error=slurm_logs/%x_%A_%a.err
#SBATCH --time=08:00:00
#SBATCH --partition=gpu
#SBATCH --gres=gpu:1
#SBATCH --cpus-per-task=8
#SBATCH --mem=32G
#SBATCH --array=0-4

set -euo pipefail
cd ${SLURM_SUBMIT_DIR}

SEEDS=(0 1 2 3 4)
SEED=${SEEDS[$SLURM_ARRAY_TASK_ID]}

python train.py \
  --dataset rsna_trauma \
  --model resnet50 \
  --lr 1e-4 \
  --seed ${SEED} \
  --amp \
  --out_dir runs \
  --run_name "rsna_resnet50_seed${SEED}_job${SLURM_JOB_ID}"
\end{lstlisting}

\subsection{Hydra + SLURM (best practice)}
Use SLURM arrays for large sweeps, Hydra overrides for config control:
\begin{lstlisting}[language=bash]
python train.py dataset=rsna_trauma model=resnet50 optim=adamw optim.lr=1e-4 seed=${SEED}
\end{lstlisting}

\subsection{Multi-GPU training (DDP)}
\begin{lstlisting}[language=bash]
#!/bin/bash
#SBATCH --job-name=mcv_ddp
#SBATCH --output=slurm_logs/%x_%j.out
#SBATCH --error=slurm_logs/%x_%j.err
#SBATCH --time=12:00:00
#SBATCH --partition=gpu
#SBATCH --gres=gpu:2
#SBATCH --cpus-per-task=12
#SBATCH --mem=64G

set -euo pipefail
cd ${SLURM_SUBMIT_DIR}

export OMP_NUM_THREADS=${SLURM_CPUS_PER_TASK}

torchrun --nproc_per_node=2 train.py \
  --dataset rsna_trauma \
  --model resnet50 \
  --lr 1e-4 \
  --seed 0 \
  --amp \
  --out_dir runs \
  --run_name "rsna_ddp2_${SLURM_JOB_ID}"
\end{lstlisting}

\section{Checklists}
\subsection{Checklist: experiment readiness}
\begin{itemize}
    \item Runs end-to-end without code edits
    \item Outputs go to \texttt{runs/<run\_name>/}
    \item Saves resolved config + cmd
    \item Logs seed, dataset, model, learning rate
    \item Saves checkpoints (best + last)
\end{itemize}

\subsection{Checklist: code review / PR checklist}
Before merging experiment code:
\begin{itemize}
    \item Argument/config surface is documented
    \item Dataset paths not hardcoded
    \item All new arguments have defaults and help messages
    \item Config groups exist (Hydra) for dataset/model/optim
    \item SLURM scripts write logs into \texttt{slurm\_logs/}
\end{itemize}

\section{Recommended Extensions (Optional)}
\begin{itemize}
    \item Use W\&B or MLflow and log full config per run
    \item Save git commit hash in run directory
    \item Save environment info in \texttt{env.txt}
    \item Add automated experiment summary script (table of metrics across runs)
\end{itemize}

\end{document}
